\scanpage{82}
% archive foliation: 81
Inversement, soit $\unsure{\mathcal P|P}_1$ (\deleted{$\unsure{\otimes}$} catégorie $\otimes$ additive) \inserted{\unsure{avec} $\otimes$ \unsure{ass.}, avec \uline{\unsure{unité}}}, munie d'un foncteur de $(\unsure{\mathcal P|P}_1,+)$ dans $(\unsure{\mathcal P|P}_1^{\mathrm{grad.}}_{\unsure{+}},\otimes)$, \deleted{\unsure{donc additif}}
\[\Gamma^*:(\unsure{\mathcal P|P}_1,\oplus)\longrightarrow(\unsure{\mathcal P|P}_1^{\mathrm{grad}},\otimes)\]
\unsure{donc}
\[\Gamma^*(0)=1\quad\text{i.e.}\quad\Gamma^{\unsure{i|\nu}}(0)=\begin{cases}1&\text{si }\nu=0\\0&\text{si }\nu>0\end{cases}\]
\inserted{\unsure{aux degr. }$0,1$}
\[\Gamma^*(X\oplus Y)\simeq\Gamma^*(X)\otimes\Gamma^*(Y)\]
(isom \unsure{gradué}, bifonctoriel en $X,Y$, compatible aux contraintes d'unité, d'ass., de comm.)
avec
\[\Gamma^i(M)=\begin{cases}0&\text{si }i<0\\1&\text{si }i=0\\M&\text{si }i=1\end{cases}\]
(isom de foncteurs en $M$). Supposons de plus donné des hom de foncteurs
\[\Gamma^{pq}\xrightarrow{\alpha_{pq}}\Gamma^p\circ\Gamma^q\]
\[(\text{i.e.}\qquad\Gamma^{pq}(X)\xrightarrow{\alpha^X_{pq}}\Gamma^p(\Gamma^q(X))\qquad)\]
satisfaisant la condition d'associativité ci-dessus, plus les conditions d'unité pour \deleted{$p,q$} l'un des deux $p,q$ égaux à $0$ ou $1$ :
\diagram{\begin{tikzcd}[column sep=large,row sep=small]
\text{(\unsure{u_1})}&\Gamma^q\arrow[r,"\alpha_{1,q}=\mathrm{id}"]\arrow[rd,bend right=20,"\alpha_{q,1}=\mathrm{id}"']&\Gamma^1\circ\Gamma^q=\mathrm{id}\circ\Gamma^q=\Gamma^q\\
&&\Gamma^q\circ\Gamma^1=\Gamma^{\unsure{1|q}}\circ\mathrm{id}=\Gamma^q
\end{tikzcd}}
\diagram{\begin{tikzcd}[column sep=large,row sep=small]
\text{(\unsure{u_0})}&\Gamma^0\arrow[r,"\alpha_{0,q}=\mathrm{id}"]\arrow[rd,bend right=20,"{\alpha_{q,0}\text{ un }\uuline{\text{iso}}}"']&\Gamma^0\circ\Gamma^q=\Gamma^0\\
&&\Gamma^q\circ\Gamma^0.
\end{tikzcd}}

